\section{Rearranging the boundary potential}\label{sec:rearrangement}

The boundary term becomes linear in the variable \(U=\Theta(v)\).
After zero extension this potential is a \(BV\) function, and may have
a boundary jump. Rearrangement must therefore retain singular variation.
We recall the formulas used below, with
\(\Per(E)=|D\mathbf1_E|(\R^n)\).
For nonnegative \(V\in BV(\R^n)\) of bounded support, coarea gives
\[
 |DV|(\R^n)=\int_0^\infty\Per(\{V>t\})\dd t;
\]
the converse holds when the right side is finite and \(V\in L^1\).
We use its Borel-set version as well
\cite[Theorem~1.19]{Gorny}; see also
\cite[Theorem~2.6]{LV}.
Euclidean isoperimetry in this normalization is
\[
 \Per(E)\ge n\omega_n^{1/n}|E|^{(n-1)/n}
\]
for bounded measurable sets, with equality for balls
\cite[Theorem~2.1]{AGNT}.

For \(f\in W^{1,1}(\Omega)\), Gauss--Green gives the zero-extension
identity
\begin{equation}\label{eq:zeroextension}
 D(f\mathbf1_\Omega)=\nabla f\,\dd x|_\Omega-
                     \Tr f\,\nu\,\mathcal H^{n-1}|_{\partial\Omega},
 \qquad
 |D(f\mathbf1_\Omega)|(\R^n)=\int_\Omega|\nabla f|+
                                          \int_{\partial\Omega}|\Tr f|.
\end{equation}
The volume and boundary measures are mutually singular.
The absolute-mean characterization also shows that the trace commutes
with a Lipschitz scalar composition.

For a Borel \(p:[0,Y]\to[0,1]\), let \(B(t)=\int_0^t p(a)\dd a\).
This is a nondecreasing 1-Lipschitz function. The Sobolev chain rule gives
\(\nabla B(U)=p(U)\nabla U\) inside \(\Omega\); the \(C^1\) version is
\cite[Proposition~3.21(2)]{Hunter}, and we spell out the representative
issue. Sobolev functions have absolutely continuous representatives on
almost every coordinate line: product tests and Fubini give the
one-dimensional weak derivative, whose integral differs from the
function by a constant \cite[Proposition~3.6]{Hunter}.
For an absolutely continuous real \(f\) and a real null Borel set \(N\),
\(f'=0\) a.e. on \(f^{-1}(N)\). Indeed, points with nonzero derivative
can be grouped into countably many sets on which the inverse
restriction of \(f\) is Lipschitz: group by integers \(k,l\) for which
\(|f(y)-f(x)|\ge |y-x|/k\) when \(0<|y-x|<1/l\), and then by intervals
of length less than \(1/l\). Their points over \(N\) are null, since
a Lipschitz map sends null sets to null sets. Apply this to the
exceptional levels of \(B'\). The ordinary chain rule holds elsewhere,
and the Lipschitz bound gives derivative zero where \(f'=0\).
Fubini proves the asserted weak chain rule with the chosen Borel \(p\).

We will also use the weighted coarea identity
\begin{equation}\label{eq:weightedcoarea}
 |D(B(V))|(\R^n)=\int_0^Y p(t)\Per(\{V>t\})\dd t
\end{equation}
for \(0\le V\le Y\) in \(BV\) of bounded support. To include flat pieces
of \(B\), put
\(\tau(q)=\sup\{t\in[0,Y]:B(t)\le q\}\) for \(0<q<B(Y)\).
Then \(\{B(V)>q\}=\{V>\tau(q)\}\). The one-dimensional area formula
\cite[Theorem~2.8]{LV}, first on indicator functions and then on
nonnegative measurable weights, changes \(dq\) to \(p(t)\dd t\).
Non-singleton fibers of \(B\) have only countably many distinct values.
This proves \eqref{eq:weightedcoarea}; its finite right side first
justifies \(B(V)\in BV\) by the coarea converse. If \(B(Y)=0\), both
sides are zero.

The scalar cost of an absolutely continuous potential drop will be
\begin{equation}\label{eq:psi}
 \psi(y,z)=\sup_{0\le p\le1}\{pz-G(y)p^2\},\qquad0\le y\le Y,\quad z\ge0.
\end{equation}
For \(G(y)>0\) this is
\[
 \psi(y,z)=
 \begin{cases}
 z^2/(4G(y)),&z\le2G(y),\\
 z-G(y),&z\ge2G(y).
 \end{cases}
\]
At \(y=0\) it is \(z\). The supremum formula shows that it is convex,
nondecreasing and 1-Lipschitz in \(z\), that \(0\le\psi(y,z)\le z\),
and that its recession slope is 1. Writing
\(\psi_g(z)=\sup_{0\le p\le1}(pz-gp^2)\), we have
\begin{equation}\label{eq:psilip}
 |\psi_{g_1}(z)-\psi_{g_2}(z)|\le |g_1-g_2| .
\end{equation}

Let \(w\) be a nonincreasing profile on \((1,R)\), taking values in
\([0,Y]\). Augment it by \(Y\) to the left of 1 and by 0 to the right
of \(R\). Its positive drop measure \(\mu=-Dw\) on \([1,R]\) has total
mass \(Y\), including the atoms \(Y-w(1+)\) at 1 and \(w(R-)\) at \(R\).
Decompose it as
\(\mu=z(r)\dd r+\mu_s\), with \(\mu_s\perp\dd r\), and define
\begin{equation}\label{eq:F}
 F(w)=\int_1^R A(r)\psi(w(r),z(r))\dd r+
                              \int_{[1,R]}A(r)\dd\mu_s(r).
\end{equation}
The singular coefficient is the recession slope of \(\psi\).
It includes jump and Cantor variation.

\begin{proposition}\label{prop:lower}
For every original admissible temperature \(v\), there is such a
profile \(w\) for which
\begin{equation}\label{eq:lower}
 \mathcal E_\Omega(v)\ge S F(w).
\end{equation}
\end{proposition}

\begin{proof}
Extend \(U=\Theta(v)\) by zero outside \(\Omega\). The \(C^1\) chain
rule and \eqref{eq:zeroextension} put \(U\) in \(BV(\R^n)\).
Completing a square gives
\[
 |\nabla v|^2\ge p(U)|\nabla U|-G(U)p(U)^2 .
\]
Since \(B(\Theta(\Tr v))\le\Theta(\Tr v)\), we obtain
\begin{equation}\label{eq:dualoriginal}
 \mathcal E_\Omega(v)\ge |D(B(U))|(\R^n)-
                                       \int_{\R^n}G(U)p(U)^2\dd x.
\end{equation}
Let \(U^*(x)=w(|x|)\) be the decreasing spherical rearrangement.
Explicitly, for \(m(t)=|\{U>t\}|\), take
\[
 w(r)=\sup\{t\in[0,Y]:m(t)>\omega_n r^n\},
\]
with empty supremum 0. Right continuity of \(m\) shows that
\(\{U^*>t\}\) is the ball of volume \(m(t)\).
Isoperimetry and the coarea converse give \(U^*\in BV\).
By \eqref{eq:weightedcoarea}, the variation of \(B(U^*)\) does not
exceed that of \(B(U)\); equimeasurability preserves the bulk term
in \eqref{eq:dualoriginal}, since its integrand vanishes at zero.
Moreover \(U=Y\) on \(K\) and its support has volume at most
\(\omega_nR^n\). Thus \(w=Y\) on \(r<1\) and \(w=0\) on \(r>R\), and
\begin{equation}\label{eq:radialdual}
 \mathcal E_\Omega(v)\ge
 \sup_p\left\{|D(B(U^*))|-\int_{\R^n}G(U^*)p(U^*)^2\dd x\right\}.
\end{equation}

We compute this supremum without assuming an absolutely continuous
inverse distribution. Complete the monotone graph of \(w\) by filling
all vertical jumps, including the endpoint jumps. From \((1,Y)\)
to \((R,0)\), use the parameter
\[
 s=r+Y-t,\qquad1\le s\le R+Y.
\]
An increase in \(s\) is exactly the radial increase plus the potential
drop. Hence its coordinate maps \(r(s),t(s)\) are 1-Lipschitz,
and, almost everywhere,
\[
 a=r'(s)\ge0,\qquad b=-t'(s)\ge0,\qquad a+b=1.
\]
There are no missing parameters: the jumps of the strictly increasing
map \(r\mapsto r+Y-w(r)\) are precisely the inserted vertical segments,
and its generalized inverse supplies the rest of the graph.
The one-dimensional area formula gives
\begin{equation}\label{eq:pushforwards}
 r_\#(a\,\dd s)=\dd r,\qquad
 t_\#(b\,\dd s)=\dd t|_{(0,Y)},\qquad r_\#(b\,\dd s)=\mu.
\end{equation}
For the last equality, a radial interval traverses exactly the drop
of \(w\) on it, with its endpoint vertical segments. The area formula
for \(t\) measures that drop. Equality on intervals determines the
finite measures, and nonnegative simple functions give all weights.

On \(a>0\), the radial inverse is unique except at countably many jump
radii; the drop density is \(b/a\). The image of \(a=0\) under \(r\)
is Lebesgue-null. Uniqueness of the Lebesgue decomposition therefore
identifies these as the absolutely continuous and singular parts of
\(\mu\), respectively. At a nonjump radius \(t(s)=w(r(s))\); a vertical
fiber has \(a=0\). Consequently
\begin{equation}\label{eq:graphcost}
 F(w)=\int_1^{R+Y}A(r(s))
       \sup_{0\le p\le1}\{p b(s)-G(t(s))p^2a(s)\}\dd s.
\end{equation}
For \(a>0\) the integrand is \(a\psi(t,b/a)\), and for \(a=0\) it
is \(b\), retaining the whole singular cost.

The radial dual in \eqref{eq:radialdual}, divided by \(S\), equals
\begin{equation}\label{eq:graphdual}
 \int_1^{R+Y}A(r(s))
       [p(t(s))b(s)-G(t(s))p(t(s))^2a(s)]\dd s
                                  -\frac{G(Y)p(Y)^2}{n}.
\end{equation}
The last term is the contribution of the core \(B_1\).
Indeed \eqref{eq:weightedcoarea} integrates the spherical crossing
radius against \(p(t)\dd t\), while polar volume is
\(S A(r)\dd r\); \eqref{eq:pushforwards} gives the two graph integrals.

It remains to choose a single Borel function of the level which
maximizes them. Choose Borel representatives of \(a,b\). A null
Borel set \(N\subset[0,Y]\) can contain \(0,Y\), the image under \(t\)
of \(b=0\), and the image of the parameter-null set where derivatives
are undefined. The area formula and the Lipschitz bound make these
images null. Outside \(N\), \(t\) has a unique inverse parameter
\(s(t)\); otherwise its fiber would contain a horizontal interval.
This generalized inverse is Borel. Set \(p=0\) on \(N\) and elsewhere
choose
\[
 p(t)=\min\{b(s(t))/(2G(t)a(s(t))),1\}\quad\text{if }a(s(t))>0,
 \qquad p(t)=1\quad\text{if }a(s(t))=0.
\]
Here \(G(t)>0\). On \(b=0\), zero is the maximizer. The part of
\(b>0\) with \(t(s)\in N\) has zero parameter measure, since its
strictly positive \(b\)-weighted integral is zero by
\eqref{eq:pushforwards}. Everywhere else the choice is a pointwise
maximizer. It also has \(p(Y)=0\), eliminating the core term.
Conversely every \(p\) gives at most the integral in
\eqref{eq:graphcost}, with a nonpositive core term.
Thus the supremum is exactly \(S F(w)\), proving \eqref{eq:lower}.
\end{proof}
